\chapter{Force fields and new parameters}
\label{ch:forcefields}

A simulation program can only move atoms it has numbers for. Those numbers
come from a force field. Standard amino acids, water and common ions are
already covered; a drug, a modified amino acid or a metal site may not be.
This chapter shows how to check what is already covered with the
\menupath{2b}{Force Field Explorer}, how far the
\menupath{2c}{Force Field Parameterizer} takes you when something is missing,
and how to keep finished parameters in your personal library.

\begin{glance}
\glanceitem{What it does}{2b shows what a force field contains and checks
  every atom of your structure against it. 2c makes new parameters for a
  ligand, a modified amino acid or a metal site. 2c option~2 stores
  parameters in your library.}
\glanceitem{You need it when}{Your structure contains anything other than
  standard amino acids, water and simple ions, and no ready-made set
  covers it.}
\glanceitem{Before you start}{For 2b, nothing (a loaded structure for the
  atom check). For 2c, a loaded structure (\key{1a}) and, for every route that
  makes charges, the Gaussian program.}
\glanceitem{Where to find it}{\key{2b} Force Field Explorer and \key{2c}
  Force Field Parameterizer in the main menu.}
\glanceitem{You get}{An atom-typing report; Gaussian input files (\file{.gjf});
  after Gaussian, \file{.frcmod} and \file{.lib} files that the Topology
  Generator (\key{5a}) can use.}
\end{glance}

\section{Force fields and parameters in plain words}
\label{sec:forcefields-words}

A \term{force field} is a set of rules and numbers that describes how atoms
attract, repel and hold together. The numbers are called
\term{parameters}. In Amber force fields each atom gets an \term{atom type}
(a short code such as \texttt{CX} for a protein C$\alpha$ or \texttt{ca} for an
aromatic carbon in GAFF2), a \term{partial charge} (a fraction of an electron
charge), and bond, angle and twist (\term{dihedral}) parameters that depend on
the atom types.

Amber keeps these numbers in a few kinds of files:

\begin{tabularx}{\linewidth}{@{}lL@{}}
\toprule
\textbf{File} & \textbf{What it holds} \\
\midrule
\file{leaprc.*} & A short script that loads a whole force field, e.g.
  \file{leaprc.protein.ff19SB}. \\
\file{.lib} or \file{.off} & A \term{residue template}: the atoms of one
  residue with their atom types and charges. \\
\file{.frcmod} & Extra or changed parameters (bonds, angles, dihedrals,
  van der Waals) that the main force field lacks. \\
\file{.mol2}, \file{.prep} & Other ways to describe a molecule with types and
  charges. \\
\bottomrule
\end{tabularx}

\medskip
For proteins the usual choice is \textbf{ff19SB} with the \textbf{OPC} water
model; ProPrep marks both with a star. For drug-like molecules Amber uses the
general force field \term{GAFF2}. A new molecule needs charges as well as
types. The careful method is \term{RESP}: charges fitted to a quantum-chemical
calculation, which needs a program such as Gaussian. A faster method is
\term{AM1-BCC}, which AmberTools can compute on its own.

\section{The beginner path}
\label{sec:forcefields-path}

Making parameters is slow and needs outside software. Work in this order and
stop as soon as everything is covered:

\begin{steps}
  \item \textbf{Check what is already covered.} Load your structure, open
        \key{2b}, load ff19SB + OPC (plus a matching add-on, such as
        \prompt{Phosphorylated AA (ff19SB)}), and run option \key{7}
        \prompt{Type Structure Atoms}. Every atom that is typed needs no new
        parameters (Section~\ref{sec:forcefields-check}).
  \item \textbf{Look for a ready-made set.} Category 12 of the Explorer lists
        69 built-in sets for hemes, iron--sulfur clusters, flavins, NAD(P),
        biopterin and Zn(Cys)$_4$ (Section~\ref{sec:forcefields-cofactors}).
        These are used through the Redox Site Preparer
        (Chapter~\ref{ch:redox}).
  \item \textbf{Only then use 2c} for what is left. Expect to stop at the
        Gaussian step and continue later (Section~\ref{sec:forcefields-2c}).
        Without Gaussian, use the AM1-BCC route in
        Section~\ref{sec:forcefields-am1bcc}.
  \item \textbf{Store finished parameters} in your library with \key{2c}
        $\to$ \key{2}, so that later projects can reuse them
        (Section~\ref{sec:forcefields-library}).
\end{steps}

\begin{watchout}
ProPrep reads the force-field files of the AmberTools it finds through the
\texttt{AMBERHOME} setting. In one of our sessions \texttt{AMBERHOME} pointed
to a different AmberTools installation, and the Explorer and tLEaP used that
one; the lists of force fields differed slightly. Always start ProPrep from
the activated ProPrep environment (Chapter~\ref{ch:installation}) and do not
set \texttt{AMBERHOME} by hand.
\end{watchout}

\section{Step by step: check what your structure needs (2b)}
\label{sec:forcefields-check}

The example is \pdb{3AJ4}, a protein with phosphoserine (SEP) and ethylene
glycol (EDO). It was found with 1a's ligand search for \typed{SEP}.

\begin{steps}
  \item Load the structure with \key{1a}. For 3AJ4 we answered \key{y} to
        \prompt{Remove hydrogen atoms?}.
  \item Type \key{2b}. The first screen lists force-field categories:
\begin{screen}
         Step 1: Select Force Field Categories
┏━━━━━━━┳━━━━━━━━━━━━━━━━━━━━━━━━━━━━┳━━━━━━━━━━━━━━━━━┓
┃ Index ┃ Category                   ┃ Available Files ┃
┡━━━━━━━╇━━━━━━━━━━━━━━━━━━━━━━━━━━━━╇━━━━━━━━━━━━━━━━━┩
│ 1     │ Protein                    │               9 │
│ 2     │ Modified Amino Acids       │               8 │
...
│ 5     │ Water                      │              12 │
│ 6     │ GAFF (General)             │               2 │
...
│ 12    │ Redox Cofactors (built-in) │              69 │
│ 13    │ Custom/User-Generated      │      User files │
└───────┴────────────────────────────┴─────────────────┘
...
Select categories (comma-separated) (1,5): §ans«1,2,5»
\end{screen}
        Categories 1--11 are Amber's own force fields. 12 holds ProPrep's
        cofactor sets and your library; 13 loads your own files.
  \item The second screen lists every force field in those categories. A
        star marks the recommended one; a warning sign marks an add-on that
        must be combined with a protein force field:
\begin{screen}
Protein
...
    9 ★ ff19SB                     Latest AMBER protein FF - use with OPC water

Modified Amino Acids  (add-ons — combine with a base protein FF)
...
   14 ⚠ Phosphorylated AA (ff19SB) Phosphorylated Ser/Thr/Tyr/His for ff19SB
...
Water
   21 ★ OPC                        4-point, optimal for ff19SB
...
   26   TIP3P                      3-point, traditional (not recommended)
...
Select force field(s) (comma-separated) (5,26): §ans«9,14,21»

Loading: leaprc.protein.ff19SB, leaprc.phosaa19SB, leaprc.water.opc
\end{screen}
  \item The Explorer menu appears. Type \key{7}
        (\prompt{Type Structure Atoms (atom typing report)}) and answer
        \key{y} to \prompt{Use this structure?}.
\begin{screen}
=================================================================
  Atom Typing Report: 3AJ4_noH.pdb
=================================================================
  Total atoms             2155
  Typed                   2142 (99.4%)
  Untyped                 13 (0.6%)
  Unique atom types       27
  Total charge (typed)    -234.2566

Untyped residues:
  A:EDO112 (4 atoms)
  A:SEP111 (1 atoms)
  B:EDO111 (4 atoms)
  B:EDO112 (4 atoms)
\end{screen}
  \item A per-residue table follows; type \key{q} to leave it. Then
        \prompt{View [a]ll atoms, [t]yped only, [u]ntyped only, [e]xport report,
        [d]one}. Here \key{u} lists the untyped atoms, \key{e} saves the
        report as a text file (\file{3AJ4_noH_atom_typing.txt}) and \key{d}
        finishes.
\end{steps}

The result is clear: the phosphoserine is already covered by the
\prompt{Phosphorylated AA (ff19SB)} add-on. Only its OXT atom was untyped,
because in this crystal SEP is the last residue of the chain. The three EDO
molecules are not covered. They are crystallisation additives, so the usual
choice is to remove them in the PDB Filter (Chapter~\ref{ch:cleaning}) rather
than to parameterize them.

\begin{watchout}
When you choose categories 1 and 5, the suggested default selection is
ff14SB + TIP3P, although the same screen calls TIP3P ``not recommended'' and
stars ff19SB and OPC. Type the numbers you want. Numbers out of range or text
are silently ignored.
\end{watchout}

\begin{tip}
The charge sums in the report are not whole numbers for a structure without
hydrogens. ProPrep explains this on screen. They only add up to the formal
residue charge when all atoms, including hydrogens, are present.
\end{tip}

\section{Exploring a force field}
\label{sec:forcefields-explore}

The Explorer menu offers:

\begin{screen}
===== Force Field Explorer =====
Loaded: leaprc.protein.ff19SB, leaprc.water.opc, leaprc.gaff2

  1. Overview (statistics & loaded files)
  2. Browse Atom Types
  3. Browse Residues
  4. Search Parameters (bond/angle/dihedral/VDW/mass)
  5. Reverse Lookup (all params for an atom type)
  6. Load Different Force Field
  7. Type Structure Atoms (atom typing report)
  b. Back
\end{screen}

\begin{itemize}
  \item \textbf{Overview} counts what was loaded (for ff19SB + OPC + GAFF2:
        1536 atom definitions, 182 residues) and lists the files.
  \item \textbf{Browse Atom Types} takes a filter such as \typed{CX}. The
        filter matches any part of the name and ignores case, so \typed{CX}
        finds both the protein type \texttt{CX} and the GAFF type
        \texttt{cx}.
  \item \textbf{Browse Residues} takes a filter such as \typed{HI}; pick a
        number to see the template, e.g. \texttt{HID} with 17 atoms, their
        types and charges.
  \item \textbf{Search Parameters}: \key{a} bond, \key{b} angle, \key{c}
        dihedral, \key{d} van der Waals, \key{e} mass. Type the atom types
        separated by spaces, e.g. \typed{XC N} gives the bond force constant
        337.0 and length 1.449~\AA{} from \file{parm19.dat}. \texttt{X} is a
        wildcard in dihedrals. Upper and lower case are different types.
  \item \textbf{Reverse Lookup} shows every bond, angle and dihedral that
        uses one atom type.
\end{itemize}

\begin{watchout}
In version 1.23.0, browsing all atom types with ff19SB loaded fills the first
pages with numbers such as \texttt{-0.00006}. These come from the ff19SB
correction grid (CMAP) and are not atom types. Always type a filter. Also,
at a filter prompt \typed{q} does not quit: it searches for names containing
``q''. Type \typed{b} to go back.
\end{watchout}

\subsection{Ready-made cofactor sets (category 12)}
\label{sec:forcefields-cofactors}

Choose \key{6} (\prompt{Load Different Force Field}), then category
\typed{12}. ProPrep lists 69 bundled parameter sets: FAD and FMN (five redox
states each), hemes (bis-His b- and c-type, Cys-ligated b-type, His/Met
c-type, with fixed-pH, constant-pH and constant-redox variants), [4Fe-4S]
clusters (oxidized FO0--FO6, reduced FD0--FD6 and FR0--FR3), NAD and NADP,
biopterin (six states) and Zn(Cys)$_4$. The table is long, because the Notes
column wraps over many lines.

At \prompt{Select set(s) (comma/range or 'none') (none)} we typed
\typed{12,69} (oxidized bis-His b-type heme and Zn(Cys)$_4$). Browse Residues
then showed the residues \texttt{HBO}, \texttt{CYZ} (zinc-bound cysteine) and
\texttt{ZNC} (the zinc).

\begin{tip}
Search by atom \emph{type}, not atom name. In the Zn(Cys)$_4$ set the zinc's
type is \texttt{ZM} and the sulfur's \texttt{SZ}; searching for \typed{ZN}
finds nothing because ZN is only the atom name. The bond \typed{ZM SZ} gives
32.3~kcal/mol/\AA$^2$ and 2.431~\AA.
\end{tip}

Category 13 opens a file browser to load your own \file{.lib}, \file{.off},
\file{.frcmod}, \file{.mol2} or \file{.prep} files. This browser moves only by
numbers and \typed{..}; you cannot type a path. Your library entries
(Section~\ref{sec:forcefields-library}) appear in category 12, under the
heading ``Redox Cofactors (built-in)'', even when they are not cofactors.

\section{The Force Field Parameterizer (2c)}
\label{sec:forcefields-2c}

\begin{screen}
══ Force Field Parameterizer Menu ══
  1. ✓ Analyze and configure residue classifications
  2. ✓ Import existing parameters into your library
  3. ○ Parameterize residues (new or resume)
       ⚠ Need to analyze residues first
  4. ○ View parameterization status
       ⚠ Need to analyze residues first
  5. ✓ Display detailed help information
\end{screen}

\subsection{Finding what needs parameters}

Type \key{1}, then \key{1} again (\prompt{Analyze}). ProPrep scans the
structure and sorts every non-standard residue into one of three kinds:
\textbf{small molecule} (an organic ligand, made with GAFF2), \textbf{modified
amino acid} (has a peptide backbone) or \textbf{metal site}. The first time,
it downloads the chemical component dictionary from \url{files.wwpdb.org}.
Examples from our tests:

\begin{itemize}
  \item \pdb{1HSG}: \prompt{MK1 (B:902) | Small Molecule}, the drug indinavir.
  \item \pdb{3AJ4}: SEP as modified amino acid, three EDO as small molecules.
  \item \pdb{1ZNI}: three Zn metal sites; the three chloride ions came out as
        \prompt{Unknown — Auto-classification could not decide}.
\end{itemize}

The analysis sub-menu can also change a classification (option 2), change the
atom-count range for small molecules (option 3, default 2--200 atoms) and edit
the table that maps modified amino acids to their parent (option 4, e.g.
\texttt{SEP}$\to$\texttt{SER}; 148 built-in entries).

\subsection{Small molecule: indinavir (MK1, 1HSG)}
\label{sec:forcefields-small}

\begin{steps}
  \item In the 2c menu type \key{3}, then \typed{1} for MK1.
  \item \prompt{Output residue name (MK1)}: press \key{Enter}. The name must
        be 1--3 letters or digits and start with a letter; \typed{9E2} is
        refused.
  \item An 8-step checklist appears:
\begin{screen}
  STRUCTURE PREPARATION
  [ ]  1.  Extract Coordinates              Extract small molecule coordinate...
  [ ]  2.  Hydrogen Addition                Analyze and optionally add hydrog...
  [ ]  3.  Gaussian Input                   Generate Gaussian input (optimiza...
  CHARGE DERIVATION
  [ ]  4.  RESP Charges                     Process Gaussian output → RESP ch...
  BONDED PARAMETER GENERATION
  [ ]  5.  Bonded Parameters                Generate missing bonded parameter...
  [ ]  6.  AMBER Topology Files             Create AMBER topology files (tLEaP)
  [ ]  7.  Parameter Refinement (optional)  Seminario / PES scan / CREST refi...
  [ ]  8.  Force Field Integration          Register lib/frcmod for the Topol...
...
  [n]ext (1. Extract Coordinates)  |  [v]iew structure  |  [s]ave & quit
  [q]uit  |  [#] to run/re-run step  |  [#s] skip optional step
\end{screen}
        Type \key{n} to run the next step, a number to run or re-run that
        step, \typed{7s} to skip the optional step 7, \key{s} to save and
        leave, \key{q} to leave without saving.
  \item \key{n}: step 1 writes \file{mk1.pdb} (45 heavy atoms).
  \item \key{n}: step 2 notices that there are no hydrogens
        (\prompt{Low hydrogen content detected (0.0\%)}) and offers to add them
        with the program \texttt{reduce}. Answer \key{y} and press \key{Enter}
        at the four option questions (NOFLIP mode, add OH/SH hydrogens, no
        nuclear distances, do not keep bond lengths). ProPrep reports
        \prompt{Added 47 hydrogen atoms}, which is right for neutral
        indinavir. Answer \key{n} to the viewer and \typed{done} to
        \prompt{Edit H atoms?}.
  \item \key{n}: step 3 writes the Gaussian input.
\begin{screen}
Net charge for MK1 (0): §ans«0»
Spin multiplicity (1): §ans«1»
Memory allocation (10GB):
Number of processors (4):
Customize calculation keywords? [y/n] (n): §ans«n»
✓ Generated Gaussian input: mk1.gjf
\end{screen}
        The \term{net charge} is the molecule's total charge (0 for a neutral
        molecule). The \term{multiplicity} describes unpaired electrons:
        1 for an ordinary molecule with all electrons paired, 2 for a
        radical.
\end{steps}

The file \file{mk1.gjf} holds two Gaussian jobs: a geometry optimization with
frequencies (\texttt{opt freq b3lyp/6-31+G(d)}), then an electrostatic
potential at \texttt{HF/6-31G(d)} for the RESP fit. If the charge is negative,
ProPrep also offers implicit solvent (\prompt{Add implicit solvation (SCRF) to
Gaussian keywords? (y)}).

\subsubsection*{The Gaussian checkpoint}

Step 4 needs Gaussian's output, \file{mk1.log}. Without it the step fails:

\begin{screen}
✗ Step 4 failed: Gaussian output not found: mk1.log
Run Gaussian first: g16 mk1.gjf
Traceback (most recent call last):
...
\end{screen}

The lines after \texttt{Traceback} look like a crash, but ProPrep is still
running and the checklist comes back. This is the point where you leave
ProPrep:

\begin{steps}
  \item Type \key{s} (save \& quit): \prompt{State saved. You can resume later.}
  \item Run Gaussian in the folder \file{small_molecule_params_MK1}:
\begin{terminal}
g16 mk1.gjf
\end{terminal}
        When it finishes you should have \file{mk1.log} and \file{mk1.chk}.
  \item Back in ProPrep, type \key{2c}, \key{3}, \typed{1}, accept the name,
        and answer \key{y} to resume:
\begin{screen}
╭──────────────────── Resume Available ─────────────────────╮
│ Found saved workflow state                                │
│ Workflow: Small Molecule Parameterization                 │
...
  Progress: 3/8 steps completed

Resume from saved state? [y/n] (y): §ans«y»
\end{screen}
\end{steps}

Answering \key{n} at \prompt{Resume from saved state?} deletes the saved state
and starts again.

\begin{needsoftware}
Steps 4--8 were not run for this manual, because Gaussian was not installed.
According to ProPrep's own description, step 4 turns the Gaussian log into
RESP charges with \texttt{antechamber}, step 5 finds missing bond parameters
with \texttt{parmchk2}, step 6 builds \file{.lib} and topology files with
tLEaP (gas phase, water or both), step 7 optionally refines parameters
(Seminario method, which also needs \texttt{formchk mk1.chk mk1.fchk}; a
dihedral scan; or CREST, which was not installed), and step 8 registers the
files for the Topology Generator and offers to store them in your library.
\end{needsoftware}

\begin{watchout}
Do not force a later step with ``Run anyway''. Running step 5 before step 4
gave the meaningless message
\texttt{sequence item 2: expected str instance, NoneType found}. Also, 2c
option~4 (\prompt{View parameterization status}) said
\prompt{No parameterizations found in workspace} although a half-finished MK1
job was saved. The saved job is still there; reach it with option~3.
\end{watchout}

\begin{watchout}
If a ligand appears several times (three EDO in 3AJ4), each copy is listed
separately. When we chose two copies, ProPrep ran the second one in the
\emph{same} folder and offered to resume the first copy's state. Answering
\key{n} there would have deleted the first copy's work. Parameterize one copy
only; the parameters apply to all copies with the same residue name.
\end{watchout}

\subsection{Modified amino acid: phosphoserine (SEP, 3AJ4)}
\label{sec:forcefields-modaa}

\begin{tip}
Check the Explorer first. SEP is already in Amber's phosphorylated
amino-acid library (Section~\ref{sec:forcefields-check}). We ran this route
only to show how far it goes.
\end{tip}

\begin{steps}
  \item \key{2c}, \key{3}, then the SEP entry. ProPrep names the parent
        residue (SER) and asks \prompt{Use SER as parent residue? (y)}:
        \key{y}.
  \item A 10-step checklist appears: Initial Structure, Backbone Angles,
        Gaussian Input, PES Scan Results, Extract Structures, ESP Generation,
        AC File, RESP Charges, Bonded Parameters, Force Field Integration.
  \item \key{n} twice: steps 1 and 2 build a capped peptide ACE--SER--NME and
        two shapes of it, with $\alpha$-helix and $\beta$-sheet backbone
        angles.
  \item \key{n}: step 3 asks for memory, processors, keywords and title
        (press \key{Enter}), then the charge, \typed{-2} for a phosphate, and
        multiplicity \typed{1}; answer \key{n} to the side-chain scan.
        ProPrep writes \file{ser_ahelix.gjf} and \file{ser_bsheet.gjf}.
  \item \key{n}: step 4 explains which Gaussian outputs it expects
        (\file{ser_ahelix.log} and \file{ser_bsheet.log}) and asks
        \prompt{Proceed (y), Pause to run Gaussian (p), or Quit (q)? (p)}.
        Type \key{p}, then \key{s}.
\end{steps}

\begin{watchout}
The two \file{.gjf} files still contain plain \emph{serine}, with no
phosphate, but carry the charge you typed ($-2$). ProPrep warns that you must
build the modification yourself. Open each file in a molecule editor such as
GaussView, add the phosphate to OG, keep the existing atom order (the frozen
backbone angles refer to atom numbers 5, 7, 9, 16 and 18), and only then run
Gaussian. Running the files unchanged would compute a serine with two extra
electrons.
\end{watchout}

Steps 5--10 (conformer extraction, ESP, multi-conformer RESP with
\texttt{residuegen}, \texttt{parmchk2}, library files) were not reached.

\begin{needsoftware}
Every 2c route that makes charges needs \textbf{Gaussian 16} (\texttt{g16}
and \texttt{formchk}): the small-molecule route, the modified-amino-acid
route and the metal-site route. ProPrep writes the Gaussian inputs and reads
the outputs; it does not run Gaussian. There is \textbf{no AM1-BCC option}
inside the app: step 4 of the small-molecule route only accepts a Gaussian
log file. Without Gaussian, ProPrep stops at the Gaussian checkpoint. For a
ligand, the route in the next section is a working alternative.
\end{needsoftware}

\section{Without Gaussian: AM1-BCC charges for a ligand}
\label{sec:forcefields-am1bcc}

AM1-BCC charges are quicker and less careful than RESP, but widely used for
drug-like ligands. The AmberTools installed with ProPrep can make them. We
did this for MK1 and then put the result into the ProPrep library.

\begin{steps}
  \item Run 2c steps 1 and 2 for the ligand (Section~\ref{sec:forcefields-small}).
        You now have \file{mk1_H.pdb}, the ligand with hydrogens, in
        \file{small_molecule_params_MK1}.
  \item Make a new empty folder, for example \file{bcc_params}, and copy
        \file{mk1_H.pdb} into it.
  \item Open a Terminal, activate the ProPrep environment
        (Chapter~\ref{ch:installation}) and go into that folder.
  \item Compute the charges and GAFF2 atom types. \texttt{-nc 0} is the net
        charge; \texttt{-rn MK1} is the residue name. This took about 4
        minutes for the 92 atoms of MK1:
\begin{terminal}
antechamber -i mk1_H.pdb -fi pdb -o mk1_bcc.mol2 -fo mol2 \
    -c bcc -nc 0 -at gaff2 -rn MK1 -s 0 -pf y
\end{terminal}
  \item Find any missing parameters:
\begin{terminal}
parmchk2 -i mk1_bcc.mol2 -f mol2 -o mk1_bcc.frcmod -s 2
\end{terminal}
  \item Make the residue template. Create a text file \file{leap.in} with
        these lines:
\begin{terminal}
source leaprc.gaff2
MK1 = loadmol2 mk1_bcc.mol2
loadamberparams mk1_bcc.frcmod
check MK1
saveoff MK1 mk1_bcc.lib
quit
\end{terminal}
        and run it:
\begin{terminal}
tleap -f leap.in
\end{terminal}
        The end of the output should read
        \texttt{Errors = 0; Warnings = 0}.
  \item Back in ProPrep, type \key{2c}, then \key{2}
        (\prompt{Import existing parameters into your library}). In the file
        browser open \file{bcc_params}, pick \file{mk1_bcc.frcmod}, then
        \file{mk1_bcc.lib}, then type \typed{skip} for the optional
        \file{.prep} file. This browser also accepts a typed path.
  \item Answer the questions:
\begin{screen}
Category (1): §ans«1»

The library defines unit MK1; the filename says mk1_bcc. tLEaP matches
the unit name and the match is case-sensitive, so MK1 is recorded as the
residue name.
...
Library entry name (residue name) (MK1):
One-line description (MK1 (GAFF, AM1-BCC/RESP charges, user-parameterized)):
§ans«Indinavir (MK1) from 1HSG, GAFF2 + AM1-BCC charges»
Parameter set name (user_MK1): §ans«MK1_AM1BCC»
Reference(s) / citation (optional; separate multiple with ' | ') ():
\end{screen}
        Category 1 is \prompt{Small molecule / ligand}; 2 is modified amino
        acid; 3 is metal site.
  \item ProPrep confirms \prompt{Parameters saved — Saved to your library}
        and says the set \prompt{will now appear in the Force Field Explorer
        and Topology Generator in any project}.
\end{steps}

\begin{tip}
Write the real charge method in the description. The suggested text says
``AM1-BCC/RESP charges'', which does not tell you later which one you used.
\end{tip}

ProPrep then prints \prompt{Imported parameters need a transformer}. A
transformer (Chapter~\ref{ch:redox}) is what links a library entry to a site
so that the Topology Generator finds it. ProPrep's advice is: load the
structure, run the Redox Site Detector, then \key{4d} option~1
(\prompt{Create transformer}), and type \typed{save} straight away to accept a
pass-through linked to \texttt{small\_molecules/MK1}. We did not carry out
this last step or use MK1 in the Topology Generator for this manual.

\section{The parameter library}
\label{sec:forcefields-library}

The library is a personal shelf of finished parameters, kept in your home
folder. The full path is

\smallskip
\centerline{\file{~/.proprep/forcefield_params/specialized_residues/}}
\smallskip

\noindent with sub-folders \file{small_molecules}, \file{modified_aa} and
\file{metal_sites}. Each entry
has a \file{metadata.json} and its \file{.frcmod} and \file{.lib} files.

\begin{itemize}
  \item \textbf{Deposit} with 2c option~2 (any time, no structure needed), or
        at the end of a parameterizer (step ``Force Field Integration'', which
        asks \prompt{Add to your library (n)}; not reached in our tests).
  \item \textbf{Find} in the Explorer, category 12. After our import the
        count rose from 69 to 71 and the MK1 sets appeared as rows 69 and
        70. Selecting one loaded \file{mk1_bcc.frcmod} and \file{mk1_bcc.lib},
        and Browse Residues showed MK1 with 92 atoms.
  \item \textbf{Reuse} in the Topology Generator of any project, through a
        transformer (see above; not tried here).
\end{itemize}

An entry has two levels of names, which ProPrep explains in a panel
\prompt{Two names, two levels}: the \textbf{entry name} says which molecule
(one folder each); the \textbf{set name} says whose numbers. One entry can
hold several sets, and exactly one is the default.

\subsection*{Versions}

Importing a set with a name that already exists gives three choices:

\begin{screen}
Set 'MK1_AM1BCC' already exists for small_molecules/MK1 [single_state/default].
  1. Save as a new version — keep the existing set; write this one as MK1_AM1BCC_v2
  2. Overwrite the existing set — replace its files and metadata with these
  3. Cancel — leave the library untouched -- the default
How would you like to proceed? (3): §ans«1»
\end{screen}

\begin{watchout}
In version 1.23.0, ``Save as a new version'' makes the new set the default
without saying so. Both versions also point to the same file names in the
same folder, so if the new files differ, the old set's files are replaced.
Give new files new names (for example \file{mk1_bcc_v2.frcmod}) before you
import a second version. Note too that row numbers in Explorer category~12
shift when you add entries.
\end{watchout}

\section{Advanced: metal sites}
\label{sec:forcefields-metal}

\begin{watchout}
This route is for experienced users. In our tests with zinc insulin
(\pdb{1ZNI}) and rubredoxin (\pdb{1IRO}, Fe(III)Cys$_4$) it never produced
usable Gaussian inputs. Read the known problems below before you start. If a
ready-made set exists (Zn(Cys)$_4$, hemes, [4Fe-4S]), use that instead via
the Redox Site Preparer (Chapter~\ref{ch:redox}).
\end{watchout}

The metal route follows the \term{MCPB} approach (the ``bonded model''): the
metal and its ligands get their own bonds, angles and charges, fitted to
quantum calculations on a \term{small model} (for force constants) and a
\term{large model} (for charges). Other metals in the structure are treated
as free ions: removed before tLEaP and put back afterwards.

You start it like the others: \key{2c}, \key{3}, then a metal site. A 15-step
checklist follows:

\begin{tabularx}{\linewidth}{@{}lL@{}}
\toprule
\textbf{Group} & \textbf{Steps} \\
\midrule
Structure preparation & 1 Structure Filtering (optional), 2 Structure
  Completeness (optional), 3 Structure Triage \\
Component parameterization & 4 Protein: FF + Protonation, 5 Water: Model
  Selection, 6 Organic Small Molecules, 7 Organometallic Small Molecules,
  8 Metal Clusters, 9 Isolated Metal Ions \\
Structure assembly & 10 Structure Recombination (tLEaP), 11 Redox Site
  Sync \\
Metal site parameterization & 12 MCPB Atom Typing, 13 Bonded Parameters,
  14 RESP Charges, 15 Force Field Integration \\
\bottomrule
\end{tabularx}

\medskip
Steps 1 and 2 reuse the PDB Filter and Structure Fixer
(Chapter~\ref{ch:cleaning}); step 4 reuses the protonation tool
(Chapter~\ref{ch:protonation}). Step 11 runs the Redox Site Detector if you
have not done so. Step 12 asks for the metal's \term{formal charge} (e.g.
\typed{+2} for Zn, \typed{+3} for Fe(III)) and its number of unpaired
electrons (\typed{0} for Zn, \typed{5} for high-spin Fe(III)), and writes
\file{small_freq.gjf} and \file{large_resp.gjf} in
\file{site_N/models/}. Step 13 then waits for the Gaussian results (run
\texttt{g16 small\_freq.gjf}, \texttt{formchk small\_freq.chk
small\_freq.fchk} and \texttt{g16 large\_resp.gjf} in that folder).

\subsection*{What happened in our tests}

\begin{itemize}
  \item \textbf{Empty Gaussian inputs.} In both 1ZNI and 1IRO, step 12
        reported \prompt{Total residues: 0} and
        \prompt{Wrote PDB file: small.pdb (0 atoms, charge=0.00)}. The
        \file{.gjf} files contained only the route line and the charge line,
        with no atoms. They cannot be run.
  \item \textbf{tLEaP failure when sites with bonds already exist.} Running 2a
        first, as ProPrep's own tip suggests, made step 10 fail with the tLEaP message
        \texttt{Fatal Error! bond: Argument \#2 is of type String\ldots} The bond
        lines lacked chain letters, a disulfide became a bond of a residue to
        itself, and metal bonds pointed to a metal that had been removed.
        Re-running 2a without bonds did not clear it. Running 2a only at
        step 11 got past step 10.
  \item \textbf{Chloride ions.} 1ZNI's Cl$^-$ ions were classed as
        ``Unknown'', then as organic molecules. Giving Amber's
        \file{atomic_ions.lib} made ProPrep take the first ion in that file
        (Ag$^+$), and tLEaP failed. Discard crystallisation ions in step 1.
  \item \textbf{Many questions.} Step 12 asks for a formal charge on every
        atom of ordinary ligand residues (10--24 prompts per site). A whole
        metal run needed 60--120 typed answers before the Gaussian step.
  \item \textbf{Suggested charges.} The suggested model charge was always 0.
        Work it out yourself: metal charge plus charged ligands, e.g.
        Fe$^{3+}$ + 4 Cys$^-$ = $-1$. The suggested multiplicity from the
        unpaired electrons was correct (5 $\to$ 6).
  \item \textbf{Protonation of ligands.} A zinc-bound His (binding through NE2)
        was suggested as HIE; HID is right. Disulfide cysteines were suggested
        as CYS. Metal-bound Cys should be CYM (ProPrep says so on screen).
  \item \textbf{Symmetry.} 1ZNI's Zn B:31 sits on a crystal symmetry axis;
        its His ligands are in neighboring copies, so 2a found
        \prompt{No nearby atoms found}. Choose a fully coordinated metal or
        build the biological assembly first.
  \item The Help screen (2c option 5) still calls the metal steps
        ``Coming soon''; the checklist above is the current one.
\end{itemize}

\section{If something goes wrong}
\label{sec:forcefields-problems}

\begin{itemize}\sloppy
  \item \textbf{The atom-type list starts with numbers like
    \texttt{-0.00006}.} Type a filter (Section~\ref{sec:forcefields-explore}).
  \item \textbf{Typing \typed{q} at a filter prompt shows odd residues.}
    Use \typed{b} to go back.
  \item \textbf{Step 4 fails with \prompt{Gaussian output not found}.} Run
    Gaussian, then resume. Without Gaussian, use the AM1-BCC route
    (Section~\ref{sec:forcefields-am1bcc}).
  \item \textbf{A Python ``Traceback'' appears after a failed step.} ProPrep
    has not crashed; read the line starting with \texttt{✗ Step} above it.
  \item \textbf{\prompt{No parameterizations found} although you saved a
    job.} Use option~3 and answer \key{y} to resume.
  \item \textbf{The Explorer or tLEaP uses the wrong AmberTools.} Activate the
    ProPrep environment and do not set \texttt{AMBERHOME} by hand.
  \item \textbf{Metal-site Gaussian inputs have no atoms.} See
    Section~\ref{sec:forcefields-metal}; there is no workaround in version
    1.23.0.
\end{itemize}

See also Chapter~\ref{ch:troubleshooting} and Chapter~\ref{ch:knownissues}.
